\documentclass[a4paper,11pt]{article}
\usepackage[utf8]{inputenc}
\usepackage[T1]{fontenc}
\usepackage[english]{babel}
\usepackage[margin=2.3cm]{geometry}
\usepackage{graphicx}
\usepackage{booktabs}
\usepackage{xcolor}
\usepackage{hyperref}
\hypersetup{colorlinks=true,linkcolor=black,urlcolor=blue}
\begin{document}
\begin{center}
  \includegraphics[width=.40\textwidth]{assets/bending-spoons-logotype.pdf}
  \par\vspace{1.2cm}
  {\LARGE Subscription Product Analytics\par}
  \vspace{.3cm}
  {\large Project Specification\par}
  \vspace{.8cm}
  {\large Yahia Ghallale\par}
  {\large Graduate Data Scientist position\par}
  \vspace{.5cm}
  {25 September 2026\par}
\end{center}
\begin{center}
\small Project demonstration prepared as part of the application. \\
All data and results in this document are synthetic and do not describe Bending Spoons.
\end{center}

\section{Problem Statement}
Subscription products generate signals across events, users, subscriptions, billing, acquisition, and experiments. Without a consistent data foundation, teams risk making decisions from incomparable metrics, immature cohorts, or forecasts that have not been evaluated on future periods.

This application demonstrates a reproducible analytics workflow that brings these sources together, defines metrics explicitly, and makes data quality visible. It helps a product team investigate engagement, conversion, churn, revenue, acquisition, and experiment impact.

\section{What the Application Does}
\begin{itemize}
  \item Generates a deterministic sample dataset for a subscription app: users, events, subscriptions, acquisition spend, and experiment assignments.
  \item Transforms CSV inputs into a three-layer DuckDB warehouse (\texttt{raw}, \texttt{clean}, \texttt{mart}), preserving the source layer and removing exact duplicate rows from the clean layer.
  \item Checks required tables, duplicate keys, required fields, impossible dates, and missing or inconsistent invoices.
  \item Presents an interactive dashboard with date, channel, country, and device filters, KPIs, time series, and cohort analysis.
  \item Compares baseline forecasts (seasonal naive, moving average, and linear trend) using temporal backtesting and MAE/MAPE metrics.
  \item Analyses an A/B test using mature conversions, confidence intervals, a statistical test, power, and minimum detectable effect (MDE).
\end{itemize}

\section{Requirements}
\subsection{Functional Requirements}
\begin{enumerate}
  \item Generate reproducible data from a configurable random seed.
  \item Build and query a local analytical warehouse.
  \item Flag anomalies that could invalidate metric interpretation.
  \item Calculate product metrics and segment the dashboard using shared filters.
  \item Compare forecasting methods with chronological backtests that avoid temporal leakage.
  \item Communicate experiment uncertainty and make recommendations proportionate to the evidence.
\end{enumerate}
\subsection{Non-functional Requirements and Constraints}
\begin{itemize}
  \item Run locally with Python 3.11+, Streamlit, and DuckDB; no cloud service is required.
  \item Keep generation, warehouse, quality, metrics, forecasting, experimentation, and presentation in separate modules.
  \item Keep the pipeline and analyses deterministic for reproducibility and review.
  \item Use synthetic data only; no personal information or real company data is included.
\end{itemize}

\section{Metrics and Assumptions}
DAU/WAU/MAU count distinct active users over rolling 1/7/30-day windows. Trial-to-paid conversion includes only trials whose full seven-day observation window has elapsed. Period churn divides effective cancellations by paid subscriptions active at the beginning of the period. Cohort retention considers a user active when they have at least one event in a month. LTV is an estimate based on monthly ARPU, an assumed 80\% gross margin, and estimated monthly churn; it is not an observed value. Acquisition spend in filtered segments is allocated in proportion to selected signups for the corresponding channel and day.

The demo simplifies monthly billing, prices, and acquisition costs. It does not model refunds, taxes, app-store fees, currencies, or reactivations. Forecasts are interpretable baselines, not promises of future performance.

\section{Architecture}
\begin{center}
\begin{tabular}{ll}
\toprule
Component & Responsibility \\
\midrule
Generator & Deterministic synthetic data and scenario consistency \\
DuckDB warehouse & Raw/clean/mart layers and analytical facts/dimensions \\
Data quality checks & Explicit validations and failure samples \\
Metrics and analysis & KPIs, retention, forecast backtests, A/B inference \\
Streamlit & Filters, charts, tables, and decision-focused summaries \\
\bottomrule
\end{tabular}
\end{center}

\section{Example Results Interpretation}
With the default seed 42, the warehouse contains 5,000 users, 665,480 clean events, and 9,276 synthetic invoices. In the onboarding experiment, among users with a mature observation window, control conversion is 59.13\% (n=2,420) and treatment conversion is 60.30\% (n=2,373): absolute lift +1.17 percentage points, 95\% CI from $-1.61$ to $+3.94$ points, $p=0.409$, estimated power 13.1\%, and MDE of about 3.94 points.

The demo recommendation is to extend the experiment before deciding: the current sample cannot reliably distinguish a useful effect from no effect. This is an illustrative recommendation based on simulated data, not a conclusion about a real product.

\section{Run Instructions}
From the project directory:
\begin{verbatim}
python3 -m venv .venv
source .venv/bin/activate
pip install -e '.[dev]'
python -m subscription_analytics.cli build --output data --seed 42 --users 5000 --days 180
streamlit run app.py
\end{verbatim}

\vfill
\noindent\footnotesize The Bending Spoons logo is included to identify the application context; it is a trademark of its respective owner. The vector artwork was obtained from the official website: \url{https://bendingspoons.com/}.
\end{document}
